\begin{appendices}
  \chapter{}
  \section{Generative AI}
  This paper uses 3 Generative AI Tools:

  \begin{enumerate}
    \item \textbf{Cursor IDE} for expanding the weakly supervised model keyword dictionary, and implementing code efficiency improvements (using GPT-5) \cite{cursor2025,openai2025chatgpt}
    \item \textbf{Elicit AI} to assist in literature exploration phase \cite{elicit}
    \item \textbf{ChatGPT} for advice, critique \& education \cite{openai2025chatgpt}
  \end{enumerate}
  None of these tools were used to generate, or otherwise automate, the content of this paper (Except in the specific case of~\ref{appendix:weak-keyword-generation})
  Underlying my use of generative AI was the imperative to avoid ``Copy \& Paste'' strategies and instead use them as research, exploration and educational tools.

  \subsection{Elicit}
  For research, Elicit is a useful tool that performs a complex literature search based on detailed prompts.
  This was valuable for exploring research \& informing the thesis more precisely.

  I would, in future, make greater use of this tool for guiding the literature review stage.

  Here's are the \emph{exact} prompts:
  \begin{prompt}{Job Description Research}
    To what extent can we use the ``Job description'' section of Job Advertisements to map, forecast or understand the demand of certain skills or job requirements? A methodology that describes time-series representation of features in the Natural Language text, with NLP or other implementations.
  \end{prompt}

  \begin{prompt}{Innovation informed AI-Skill Forecast Research}
    Can innovation signal data (research papers, patents or open source software (e.g. \gh{})) improve the strength of skill demand forecasts derived from job postings (job boards like LinkedIn or Indeed), particularly for AI related skills? E.g. ``prompt engineering'' or ``MLOps''
  \end{prompt}

  \subsection{ChatGPT}
  For Code, ChatGPT was used for advice, debugging code or learning new programming frameworks.
  For writing, it was used as a critique assistant, to help explore ideas and critically assess the coherence of arguments or passages.

  Here are some example prompts (adjusted, on a case-by-case basis):
  \begin{prompt}{Conceptual Understanding}
    Advise me on how to think about \textbf{\{Concept/Tool/Framework/Method\}} conceptually.
    I’d like to know the mental models or guiding intuitions that professionals use.
  \end{prompt}
  \begin{prompt}{Conceptual Understanding}
    Below is my summary of what I currently understand about \textbf{\{Concept/Tool/Framework/Method\}}.
    Please critique it: point out what is accurate, what is incomplete, and
    where I’ve misunderstood. The goal is to refine my own understanding,
    not to receive a polished replacement text.
  \end{prompt}

  Much of the time I use ChatGPT for workflow and tool advice, especially with regards to using Latex:

  \begin{prompt}{Workflow advice}
    I’m setting up a research workflow using \textbf{\{Tool/Method\}}.
    Can you suggest general strategies, potential pitfalls, and best practices?
    My goal is to design the workflow myself, but I’d like conceptual guidance.
  \end{prompt}

  \begin{prompt}{\LaTeX Help}
    Advise me on strategies for managing a large dissertation in LaTeX.
    For example, what are best practices for splitting chapters into separate files,
    managing bibliographies, and handling figures/tables so that the project
    remains maintainable?
  \end{prompt}

  \begin{prompt}{Code Critique}
    Below I include a piece of code.
    Rather than rewriting it for me,
    can you summarize your understanding of its purpose, critique its efficiency, point out
    any readability/maintainability issues, and suggest what *principles* I should look at to improve it?
    \textbf{\{Code\}}
  \end{prompt}

  \begin{prompt}{Writing Critique}
    I am writing a dissertation in \textbf{\{Dissertation Context\}}.
    I will paste a passage of my draft below.
    Please provide a detailed critique focusing on:
    \begin{enumerate}
      \item clarity of expression and readability,
      \item academic style and tone,
      \item logical flow and argument structure, and
      \item any redundancy, ambiguity, or opportunities for conciseness.
    \end{enumerate}
    Give only high-level feedback (structure, argument, coherence) - do not provide any specific rephrasing or line edits.
    Assume I am aiming for a professional, dissertation-quality standard.
    Here is the passage:
    \textbf{\{Text Passage (No more than a paragraph at a time)\}}
  \end{prompt}

  Each of the above prompts are used as guiding structures, and much of the time were adapted in scenario-specific follow-ups.

  \subsection{Cursor}

  \subsubsection{Code Improvements}

  \textbf{Cursor} was used to advise on code improvements that were made strictly \emph{after} a working solution was in place.
  The purpose of this was to improve both the computational efficiency and readability of the codebase.
  Further, it serves as an educational tool for improving my coding skills.

  These adjustments were not made using the ``Agent'' feature of Cursor, but with ``Ask'', such that it provides its suggestion for implementation at my discretion.

  Here's an example prompt for this use-case:

  \begin{prompt}{Code Improvement Suggestions}
    Within \textbf{\{Context File\}} in lines \textbf{\{Context Lines\}} is a piece of code. Currently it is functional.
    Rather than rewrite it for me, can you summarize your understanding of its purpose, critique its efficiency, point out
    any readability/maintainability issues, and suggest what *principles* I should look at to improve it?
  \end{prompt}

  \subsubsection{Weak-Model Keyword Generation} \label{appendix:weak-keyword-generation}
  This is sole exception case where generated content was used: Instructing Cursor (using GPT-5) to construct the keyword dictionary for the weakly-supervised model.
  Its task here was to (In iterations of $50$ items; $4$ runs $\rightarrow 200$ items):

  \begin{itemize}
    \item Expand a \emph{Seed List} of keywords (Obtained by iteratively scraping \gh{} topics)
    \item Match keywords to their commonly used acronyms
    \item Categorize the keywords
  \end{itemize}

  The grounding of this implementation in a \emph{Seed List} was \textbf{instrumental} in making this effort defensible.

  The purpose of this was to reduce manual \& laborious data entry work - it should be noted that the only practical advantage of this implemenation was to save time and effort.
  This entire process was audited and refined manually thereafter, to ensure output quality.

  %   Here is the prompt for \textbf{Keyword list curation}:
  \begin{prompt}{Keyword List Curation}
    In \textbf{\{Context File\}} you'll find a list of keywords relating to Artificial Intelligence \& related \ghrepo{} topics.
    Your task is to expand this keyword list, to a total of no more than 50.
    The resulting list should contain core concepts, tools, programming languages and domain-specific terms.
    If a keyword is commonly referred to by an acronym, these should be paired with its relevant verbose version.
    The acronym field should be left null for terms without commonly used acronyms.
    Be sure to preserve the casing for each keyword - e.g. ``pytorch'' should be ``PyTorch''.
    Make sure not to duplicate existing entries.
  \end{prompt}

  %   And for \textbf{Categorization}:
  \begin{prompt}{Categorization}
    Within \textbf{\{Context File\}} are keywords that relate to artificial intelligence (including their acronyms).
    Create a new field of categories for each keyword into a set of no more than 25 different categories.
    These categories should be informative, and concise; if deemed necessary, one term can belong to multiple categories (no more than 3).
  \end{prompt}

  \section*{Data Sources}
  \begin{table}[ht]
    \hskip-2.0cm
    \caption{Overview of datasets used for Job Postings}
    \label{table:job-data-overview}
    \small
    \centering
    \renewcommand{\arraystretch}{1.5}
    \begin{tabular}{p{3.5cm} p{3cm} p{6.5cm} p{2cm}}
        \textbf{Dataset Name}                      & \textbf{Date Range}     & \textbf{Details}                                                  & \textbf{Source}           \\
        \hline
        1000 ML Jobs US                            & 12/2023 - 04/2025       & Mostly within March \& April of 2025, all jobs are ML related     & \cite{__MachineLearninga} \\
        Linkedin Data Engineer Job Postings        & 17/12/2023 - 17/12/2023 & Extra "skills" column - This is appended to "description"         & \cite{asaniczka_2023}     \\
        Data Science Job Postings \& Skills (2024) & 19/01/2024 - 21/01/2024 & Extra "skills" column - This is appended to "description"         & \cite{asaniczka_2024}     \\
        LinkedIn Job Postings (2023 - 2024)        & 23/03/2024 - 20/04/2024 & --                                                                & \cite{arsh_koneru_2024}   \\
        Job Postings from Ireland (October 2022)   & 01/10/2022 - 31/10/2022 & Ireland only; Sample dataset from Techmap.io's commercial product & \cite{__JobPostings}      \\
        US Job Postings from 2023-05-05            & 05/05/2023 - 05/05/2023 & USA Only; Sample dataset from Techmap.io's commercial product     & \cite{__USJoba}           \\
    \end{tabular}
\end{table}

  \section*{Third Party Code and Software Libraries}\label{appendix:software-libaries}

  \begin{enumerate}
    \item Whatever Software I end up using
  \end{enumerate}

  \subsection{ESCO Skill Classification}
  This publication uses the ESCO classification of the European Commission.
\end{appendices}